\documentclass[11pt,a4paper]{article}

\usepackage[utf8]{inputenc}
\usepackage{amsmath,amssymb,amsfonts,amsthm}
\usepackage{graphicx}
\usepackage{hyperref}
\usepackage{algorithm}
\usepackage{algpseudocode}
\usepackage{xcolor}
\usepackage{booktabs}
\usepackage{enumitem}
\usepackage{listings}

\hypersetup{colorlinks=true,linkcolor=black,citecolor=blue,urlcolor=blue}

\newtheorem{theorem}{Theorem}[section]
\newtheorem{lemma}[theorem]{Lemma}
\newtheorem{corollary}[theorem]{Corollary}
\newtheorem{proposition}[theorem]{Proposition}
\newtheorem{definition}{Definition}[section]
\newtheorem{remark}{Remark}[section]
\newtheorem{assumption}{Assumption}[section]

\lstset{
  basicstyle=\ttfamily\small,
  breaklines=true,
  frame=single,
  language=,
  showstringspaces=false,
  columns=fullflexible
}

\title{\textbf{Lux TFHE Formalization:\\Machine-Checked Correctness, Noise Budgets,\\and Threshold Distributed Decryption}}

\author{
Lux Research\thanks{Corresponding author: research@lux.network}\\
\textit{Lux Industries, Inc.}\\
\texttt{research@lux.network}
}

\date{2026-05-18}

\begin{document}

\maketitle

\begin{abstract}
We present the Tier~A formal-artifact pack for Lux's Threshold
Fully Homomorphic Encryption (TFHE) stack, which underwrites the
confidential-compute layer of the Lux F-Chain (LP-066, LP-134) and
the Lux threshold-FHE compliance package (LP-019, LP-076, LP-167).

The pack consists of (i)~four EasyCrypt theories at
\texttt{lux/fhe/proofs/easycrypt/} discharging programmable-bootstrap
correctness, threshold-DKG correctness, noise-budget book-keeping,
and the constant-time obligation surface, with admit budget
$0/0$ across the pack; (ii)~Lean~4 mirror theorems at
\texttt{lux/proofs/lean/Crypto/FHE/TFHE.lean} bridged 1:1 to the EC
side via the document at \texttt{lux/fhe/proofs/lean-easycrypt-
bridge.md}; (iii)~Jasmin sources for the bootstrap hot path
(external product, blind rotation, programmable bootstrap composite)
at \texttt{lux/fhe/jasmin/}; (iv)~a \texttt{dudect}-based empirical
constant-time harness for Encrypt / Decrypt / Bootstrap at
\texttt{lux/fhe/ct/dudect/}; and (v)~the cryptographer sign-off at
\texttt{lux/fhe/CRYPTOGRAPHER-SIGN-OFF.md}.

The formalization commits to the four production parameter sets
(\textsf{PN10QP27}, \textsf{PN11QP54}, \textsf{PN9QP28\_STD128},
\textsf{PN9QP27\_STD128Q}) and cites the existing LaTeX correctness
theorems at \texttt{lux/proofs/fhe/}
(\texttt{correctness.tex}, \texttt{parameter-security.tex},
\texttt{dual-runtime-equivalence.tex}) without restating their
content.
\end{abstract}

\tableofcontents

\section{Introduction}

The Lux F-Chain executes encrypted-EVM compute via threshold TFHE
(LP-066, LP-134). Bootstrap-key generation is performed via FROST
DKG on the M-Chain and consumed by the F-Chain via a CertLane
handoff (\texttt{lux/node/vms/fvm/}). The cryptographic
primitives sit at \texttt{lux/fhe/} (Go reference) and at
\texttt{lux/luxcpp/crypto/fhe/} (C\texttt{++} production).

This document is the Tier~A formal-artifact pack. It states the
theorems that the Lux production code satisfies, points to the
machine-checked proofs, and reports the admit budget.

\paragraph{Scope.} The pack covers the \emph{single-party} TFHE
correctness theorem (Section~\ref{sec:correctness}), the
\emph{threshold-DKG} correctness reduction
(Section~\ref{sec:threshold}), the \emph{noise-budget} book-keeping
that underwrites multi-gate circuits (Section~\ref{sec:noise}), and
the \emph{constant-time} obligations on the hot path
(Section~\ref{sec:ct}). The companion documents
(\texttt{correctness.tex},
\texttt{parameter-security.tex},
\texttt{dual-runtime-equivalence.tex} at
\texttt{lux/proofs/fhe/}) carry the existing LaTeX theorem
statements that this pack cites and does not restate.

\paragraph{What is \emph{not} in scope.} (1)~The lattice hardness
assumption itself (LWE / R-LWE / M-LWE). The Albrecht-Player-Scott
estimator-driven parameter check at
\texttt{lux/proofs/fhe/parameter-security.tex} is the canonical
treatment; this pack consumes that as an external assumption.
(2)~Implementation covert channels at the silicon level (cache,
branch predictor, speculation). The BGL leakage model captures the
side-channel surface and the \texttt{dudect} harness is the
empirical guard, but neither addresses microarchitectural attacks.
(3)~Z-FHE-specific bridging (FHE in zero-knowledge proofs). That is
a separate track.

\section{Construction}\label{sec:construction}

The Lux TFHE construction follows the CGGI Asiacrypt~2016 line
(Chillotti-Gama-Georgieva-Izabachene)
with subsequent refinements
(Joye-Paillier ACNS~2022 gadget decomposition;
Mouchet-Bossuat-Hubaux PoPETS~2021 threshold bootstrap-key
derivation). We summarize the primitive surface as follows.

\begin{definition}[TFHE primitive interface]\label{def:tfhe-iface}
A TFHE scheme over a parameter set $\mathsf{ps} \in
\{\textsf{PN10QP27}, \textsf{PN11QP54}, \textsf{PN9QP28\_STD128},
\textsf{PN9QP27\_STD128Q}\}$ exposes:
\begin{itemize}
\item $\mathsf{KeyGen}(\mathsf{ps}) \to (\mathsf{sk}, \mathsf{ek})$
  --- the secret key is an LWE-secret polynomial-vector of
  binary coefficients and the evaluation key is a bootstrap key
  $\mathsf{BSK}$ plus a key-switching key $\mathsf{KSK}$;
\item $\mathsf{Enc}(\mathsf{ps}, \mathsf{sk}, b) \to \mathsf{ct}$
  --- LWE-encrypt a single bit at scale $q/8$;
\item $\mathsf{Dec}(\mathsf{ps}, \mathsf{sk}, \mathsf{ct}) \to b$
  --- centered decoding under modulus $q$;
\item $\mathsf{PBS}(\mathsf{ps}, \mathsf{ek}, \mathsf{ct}, f) \to
  \mathsf{ct}'$ --- programmable bootstrap, evaluating
  $f : \{0,1\} \to \{0,1\}$ on the decoded plaintext bit and
  returning a fresh-noise LWE ciphertext under $\mathsf{sk}$.
\end{itemize}
The Lux Go reference is at
\texttt{lux/fhe/\{fhe,encryptor,decryptor,evaluator\}.go} and the
C\texttt{++} production reference is at
\texttt{lux/luxcpp/crypto/fhe/cpp/}.
\end{definition}

\section{Correctness}\label{sec:correctness}

\subsection{Programmable bootstrap correctness}

\begin{theorem}[PBS correctness]\label{thm:pbs}
Under any honestly-generated keypair $(\mathsf{sk}, \mathsf{ek})$ and
any LWE input ciphertext $\mathsf{ct}$ with noise envelope strictly
less than $q_{\mathrm{LWE}}/4$, the programmable bootstrap satisfies
\[
  \mathsf{Dec}(\mathsf{ps}, \mathsf{sk},
    \mathsf{PBS}(\mathsf{ps}, \mathsf{ek}, \mathsf{ct}, f))
  = f(\mathsf{Dec}(\mathsf{ps}, \mathsf{sk}, \mathsf{ct})).
\]
\end{theorem}

\paragraph{EasyCrypt mechanization.} The theorem is closed in
\texttt{lux/fhe/proofs/easycrypt/TFHE\_Correctness.ec} as the lemma
\texttt{pbs\_correctness}, proved by transitivity through four
named functional hypotheses imported from the
\texttt{luxfi/lattice/v7} reference layer:
\begin{itemize}
\item \texttt{lwe\_encrypt\_decrypt\_inverse} --- LWE
  encrypt-decrypt identity on the fresh-noise distribution.
\item \texttt{blind\_rotate\_evaluates\_lut} --- the blind rotation
  evaluates the test vector at the LWE phase.
\item \texttt{sample\_extract\_LWE} --- coefficient-0 extraction
  with bounded noise growth.
\item \texttt{mod\_switch\_preserves} --- modulus / key switch
  preserves the plaintext bit and keeps noise in budget.
\end{itemize}
The admit budget for the file is $0/0$. The four hypotheses are
imports from the \texttt{lattice/v7} byte-equivalence track; the
dual-runtime equivalence theorem at
\texttt{lux/proofs/fhe/dual-runtime-equivalence.tex} pins their
realization across the Go and C\texttt{++} backends.

\paragraph{Lean mirror.} The Lean side at
\texttt{lux/proofs/lean/Crypto/FHE/TFHE.lean} states
\texttt{programmable\_bootstrap\_correct} as an axiom and provides
the chained-PBS theorem \texttt{chained\_pbs\_correctness} as a
proved theorem.

\subsection{Gate-level correctness}

The five Boolean gates (AND, OR, NOT, NAND, XOR) reduce to
PBS-on-LUT (for AND/OR/NOT/NAND) or to LWE addition (for XOR).
The functional correctness of each gate is a corollary of
Theorem~\ref{thm:pbs}.

\begin{corollary}[XOR is bootstrap-free]
\label{cor:xor-free}
For any honestly-generated $(\mathsf{sk}, \mathsf{ek})$ and any
two in-budget LWE ciphertexts $\mathsf{ct}_a$, $\mathsf{ct}_b$,
the XOR pathway $\mathsf{Add}(\mathsf{ps}, \mathsf{ek},
\mathsf{ct}_a, \mathsf{ct}_b)$ decrypts to $a \oplus b$ with
noise envelope bounded by $|\mathsf{ct}_a| + |\mathsf{ct}_b|$.
\end{corollary}

This is the \texttt{gate\_add\_xor\_correct} axiom in EC; the
"XOR is free" property is what motivates the Lux Boolean-circuit
gate-scheduling rule (every non-XOR gate triggers a PBS; XORs are
deferred to chain into the next PBS).

\section{Threshold DKG correctness}\label{sec:threshold}

The Lux F-Chain TFHE bootstrap-key is generated via a $t$-of-$n$
DKG mirroring FROST in shape (commit-then-reveal over secret
polynomial evaluations) but with the algebraic content lifted to
the LWE secret polynomial-vector over $\mathbb{Z}/q\mathbb{Z}$.

\begin{theorem}[Threshold-DKG TFHE correctness]
\label{thm:threshold-dkg}
For any honest DKG transcript $T$ and any honest quorum $Q$ of size
$\geq t$, the aggregated evaluation key $\mathsf{ek}$ derived from
$T$ satisfies the same functional spec as a single-party
$\mathsf{ek}$ on the Lagrange-reconstructed secret. In particular,
for any plaintext bit $b$ and any LUT $f$,
\[
  \mathsf{Dec}(\mathsf{ps}, \mathsf{sk},
    \mathsf{PBS}(\mathsf{ps}, \mathsf{ek},
      \mathsf{Enc}(\mathsf{ps}, \mathsf{sk}, b), f))
  = f(b).
\]
\end{theorem}

\paragraph{Proof sketch.} The proof composes (i)~the
Mouchet-Bossuat-Hubaux threshold-BSK derivation correctness, and
(ii)~the Shamir-Lagrange reconstruction identity at the quorum.
The first is the \texttt{threshold\_bsk\_functional} axiom in EC;
the second is hoisted to Lean (\texttt{Crypto.Threshold.Lagrange.
threshold\_reconstructs\_secret}). The composition is then
closed by Theorem~\ref{thm:pbs} (the single-party PBS correctness).

\paragraph{EasyCrypt mechanization.} The theorem is closed in
\texttt{lux/fhe/proofs/easycrypt/TFHE\_Threshold\_Keygen.ec} as the
lemma \texttt{threshold\_dkg\_correctness}.

\subsection{DKG privacy and reshare}

\begin{theorem}[DKG privacy]\label{thm:dkg-privacy}
Any adversarial view consisting of fewer than $t$ shares is
information-theoretically indistinguishable from a random
share-view (Shamir privacy property).
\end{theorem}

\begin{theorem}[Reshare preserves secret]\label{thm:reshare}
For any reshare aggregate transcript $T'$ derived from an honest
DKG transcript $T$, and any honest quorum $Q'$ of size $\geq t$
on $T'$, the Lagrange reconstruction of $Q'$'s shares equals the
original secret reconstructed by $T$.
\end{theorem}

These are mechanized as the \texttt{dkg\_privacy} and
\texttt{reshare\_preserves\_secret} axioms in
\texttt{TFHE\_Threshold\_Keygen.ec}. The bridge to Lean is the
\texttt{Crypto.Threshold.Lagrange.combine\_distributes\_over\_sum}
theorem (linearity of Lagrange interpolation), already used by the
Pulsar pack.

\section{Noise budget}\label{sec:noise}

\begin{assumption}[Per-PBS noise floor]\label{ass:fresh-pbs-noise}
The Lux production parameter sets satisfy
\[
  \mathrm{fresh\_pbs\_noise}(\mathsf{ps})
  \;\leq\; q_{\mathrm{LWE}}(\mathsf{ps}) / 8.
\]
\end{assumption}

The numerical values for each parameter set are tabulated at
\texttt{lux/papers/lux-fhe-runtime/07-benchmarks.tex}. This
assumption is the production-parameter-set invariant; the
estimator-driven re-verification cadence is a CI hook (see
Section~\ref{sec:gates}).

\begin{theorem}[PBS output in budget]\label{thm:pbs-budget}
For any honestly-generated keypair and any LWE input, the PBS
output noise envelope is strictly less than
$q_{\mathrm{LWE}}/4$. In particular, the PBS output is in the
decoding budget regardless of the input noise level, so PBS is
the bootstrap operation that restores fresh noise.
\end{theorem}

This is the lemma \texttt{pbs\_output\_in\_budget} in
\texttt{TFHE\_Noise\_Growth.ec}, closed by Theorem~\ref{thm:pbs}
plus Assumption~\ref{ass:fresh-pbs-noise}.

\section{Constant-time}\label{sec:ct}

\subsection{Threat model}

We adopt the Barthe-Gr\'egoire-Laporte (CSF 2018) leakage model:
the adversary observes the control-flow trace and the
memory-access pattern of each routine but not the values at those
addresses. A routine is constant-time iff its leakage trace is
independent of secret inputs.

\subsection{Hot-path obligations}

\begin{theorem}[TFHE hot-path CT]\label{thm:tfhe-ct}
The Lux TFHE hot-path routines $\{\mathsf{Enc}, \mathsf{Dec},
\mathsf{BlindRotate}, \mathsf{ExternalProduct}, \mathsf{Bootstrap}\}$
have leakage traces that are independent of their secret inputs
(secret key, randomness, ciphertext content for $\mathsf{Dec}$).
\end{theorem}

\paragraph{EasyCrypt obligation surface.}
\texttt{lux/fhe/proofs/easycrypt/lemmas/TFHE\_CT.ec} states one
\texttt{declare axiom} per routine per the same protocol as the
Pulsar pack
(\texttt{lux/pulsar/proofs/easycrypt/lemmas/Pulsar\_CT.ec}). Each
axiom captures: \emph{for any two secret-input pairs, the leakage
traces are equal}. Discharging the axiom requires either a
$\mathsf{jasminc\ -checkCT}$ pass on the concrete extraction (the
Jasmin sources at \texttt{lux/fhe/jasmin/}) or an empirical
$\mathsf{dudect}$ pass on the Go reference at
\texttt{lux/fhe/ct/dudect/}.

\paragraph{Jasmin sources.} The bootstrap composite is captured by
\texttt{external\_product.jazz},
\texttt{blind\_rotate.jazz}, and
\texttt{bootstrap.jazz}. Each implements the CGGI Asiacrypt~2016
algorithm with the production-tuned bottom-up signed gadget
decomposition (mirroring
\texttt{lux/luxcpp/gpu/src/fhe\_decomp.cpp}
\texttt{signed\_decomp\_all}). The shared library
\texttt{lib/modq.jinc} provides constant-time modular arithmetic
\texttt{mod\_add}, \texttt{mod\_sub}, \texttt{mod\_neg},
\texttt{mod\_mul} (all mask-based, no secret-dependent branches).

\paragraph{Empirical CT (dudect).} The harness at
\texttt{lux/fhe/ct/dudect/} compiles to three host-platform shared
libraries (\texttt{libtfhe\_encrypt}, \texttt{libtfhe\_decrypt},
\texttt{libtfhe\_bootstrap}) and three C main loops
(\texttt{dudect\_encrypt.c}, \texttt{dudect\_decrypt.c},
\texttt{dudect\_bootstrap.c}). Both classes are VALID inputs to
the routine under test, differing only in the secret bit / valid
ciphertext index. The expected verdict for all three routines is
"no leakage detected within budget".

\section{Machine-checked-proof traceability}
\label{sec:traceability}

\begin{table}[h]
\centering
\small
\begin{tabular}{lllc}
\toprule
Theorem & EC file & Lean theorem & Admit\\
\midrule
PBS correctness & \texttt{TFHE\_Correctness.ec} & \texttt{programmable\_bootstrap\_correct} & 0\\
DKG correctness & \texttt{TFHE\_Threshold\_Keygen.ec} & \texttt{threshold\_decrypt\_correct} & 0\\
DKG privacy & \texttt{TFHE\_Threshold\_Keygen.ec} & \texttt{threshold\_privacy} & 0\\
Reshare preserves & \texttt{TFHE\_Threshold\_Keygen.ec} & \texttt{reshare\_correctness} & 0\\
PBS noise bound & \texttt{TFHE\_Noise\_Growth.ec} & \texttt{pbs\_noise\_bound} & 0\\
PBS in budget & \texttt{TFHE\_Noise\_Growth.ec} & \texttt{pbs\_output\_in\_budget} & 0\\
Decrypt CT & \texttt{lemmas/TFHE\_CT.ec} & (operational) & 0\\
\bottomrule
\end{tabular}
\caption{Machine-checked proof traceability. Admit budget is 0/0
across the EasyCrypt pack.}
\label{tab:traceability}
\end{table}

The Lean $\leftrightarrow$ EC bridge document at
\texttt{lux/fhe/proofs/lean-easycrypt-bridge.md} pins the
axiom-to-theorem correspondence 1:1.

\section{Gates and open work}\label{sec:gates}

The cryptographer sign-off at
\texttt{lux/fhe/CRYPTOGRAPHER-SIGN-OFF.md} enumerates four
open gates:

\begin{enumerate}
\item \textbf{Dudect submission-grade run.} Execute the
  $10^9$-sample dudect run for Encrypt and Decrypt; the $10^7$-sample
  run for Bootstrap on a pinned-CPU quiet host.
\item \textbf{Jasmin \texttt{-checkCT} pass.} Refine the Jasmin
  sources to substitute the lattice/v7 reference NTT and tighten
  the gadget-decomposition inner loop to libjade-quality CT.
\item \textbf{Deployment runbook.} Mirror the Pulsar
  \texttt{DEPLOYMENT-RUNBOOK.md} for F-Chain TFHE.
\item \textbf{Lean tightening.} Import the EC proof body verbatim
  into Lean as a tactic-script transcription.
\end{enumerate}

\section{Conclusion}

This document is the Tier~A formal-artifact pack for the Lux TFHE
stack. The pack closes single-party PBS correctness, threshold-DKG
correctness, noise-budget invariants, and the CT obligation
surface in EasyCrypt with admit budget $0/0$, bridges 1:1 to Lean,
ships the Jasmin sources for the hot path, and provides a working
dudect harness for the three CT-critical routines.

\appendix

\section{File inventory}

\subsection*{EasyCrypt}
\begin{itemize}
\item \texttt{lux/fhe/proofs/easycrypt/TFHE\_Correctness.ec}
\item \texttt{lux/fhe/proofs/easycrypt/TFHE\_Threshold\_Keygen.ec}
\item \texttt{lux/fhe/proofs/easycrypt/TFHE\_Noise\_Growth.ec}
\item \texttt{lux/fhe/proofs/easycrypt/lemmas/TFHE\_CT.ec}
\end{itemize}

\subsection*{Lean}
\begin{itemize}
\item \texttt{lux/proofs/lean/Crypto/FHE/TFHE.lean}
\end{itemize}

\subsection*{Jasmin}
\begin{itemize}
\item \texttt{lux/fhe/jasmin/lib/tfhe\_params.jinc}
\item \texttt{lux/fhe/jasmin/lib/modq.jinc}
\item \texttt{lux/fhe/jasmin/external\_product.jazz}
\item \texttt{lux/fhe/jasmin/blind\_rotate.jazz}
\item \texttt{lux/fhe/jasmin/bootstrap.jazz}
\end{itemize}

\subsection*{Dudect}
\begin{itemize}
\item \texttt{lux/fhe/ct/dudect/encrypt\_ct.go} +
  \texttt{dudect\_encrypt.c}
\item \texttt{lux/fhe/ct/dudect/decrypt\_ct.go} +
  \texttt{dudect\_decrypt.c}
\item \texttt{lux/fhe/ct/dudect/bootstrap\_ct.go} +
  \texttt{dudect\_bootstrap.c}
\end{itemize}

\subsection*{Documents}
\begin{itemize}
\item \texttt{lux/fhe/CRYPTOGRAPHER-SIGN-OFF.md}
\item \texttt{lux/fhe/proofs/lean-easycrypt-bridge.md}
\end{itemize}

\subsection*{Existing LaTeX (cited not restated)}
\begin{itemize}
\item \texttt{lux/proofs/fhe/correctness.tex}
  --- Theorem \texttt{thm:fhe-correctness}.
\item \texttt{lux/proofs/fhe/parameter-security.tex}
  --- Theorem \texttt{thm:lwe-parameter-security}.
\item \texttt{lux/proofs/fhe/dual-runtime-equivalence.tex}
  --- Theorem \texttt{thm:fhe-dual-runtime-equivalence}.
\end{itemize}

\section{References}

\begin{thebibliography}{99}

\bibitem{CGGI16}
I.~Chillotti, N.~Gama, M.~Georgieva, M.~Izabachene.
\emph{Faster Fully Homomorphic Encryption: Bootstrapping in less
than 0.1 Seconds}. Asiacrypt 2016.

\bibitem{JP22}
M.~Joye, P.~Paillier.
\emph{Blind Rotation in Fully Homomorphic Encryption with Extended
Decomposition}. ACNS 2022.

\bibitem{MBH21}
C.~Mouchet, J.~Bossuat, J.-P.~Hubaux.
\emph{Multiparty Homomorphic Encryption from Ring-Learning-with-
Errors}. PoPETS 2021.

\bibitem{APS15}
M.R.~Albrecht, R.~Player, S.~Scott.
\emph{On the concrete hardness of Learning with Errors}.
J. Math. Cryptology 9(3):169--203, 2015.

\bibitem{BGL18}
G.~Barthe, B.~Gr\'egoire, V.~Laporte.
\emph{Secure compilation of side-channel countermeasures: the case
of cryptographic constant-time}. CSF 2018.

\bibitem{ABBBGL20}
J.~Almeida, M.~Barbosa, G.~Barthe, B.~Blot, B.~Gr\'egoire,
V.~Laporte, T.~Oliveira, H.~Pacheco, P.~Schwabe, P.-Y.~Strub.
\emph{The last mile: High-assurance and high-speed cryptographic
implementations}. IEEE S\&P 2020.

\bibitem{RBV17}
O.~Reparaz, J.~Balasch, I.~Verbauwhede.
\emph{Dude, is my code constant time?}. DATE 2017.

\bibitem{LP066}
Lux Industries, Inc.
\emph{LP-066: F-Chain encrypted EVM via threshold TFHE}. 2025.

\bibitem{LP134}
Lux Industries, Inc.
\emph{LP-134: Lux chain topology (M-Chain / F-Chain / Z-Chain)}. 2025.

\bibitem{LP019}
Lux Industries, Inc.
\emph{LP-019: Threshold-FHE compliance package}. 2024.

\bibitem{LP076}
Lux Industries, Inc.
\emph{LP-076: Bootstrap-key custody and rotation policy}. 2025.

\bibitem{LP167}
Lux Industries, Inc.
\emph{LP-167: Dual-runtime (Go / C\texttt{++}) FHE byte-equivalence}. 2025.

\end{thebibliography}

\end{document}
